\begin{textsample}{NEG Dim 1 – vem\_ed\_28 – Score 22.00 – vem\_ed\_28/t146.txt}  \label{ex:f1_neg_185}
Leitores Aos Osvaldo Succi Jr.
Coordenador PCIs No número anterior deste informativo VEm (janeiro / fevereiro 2025), sintetizamos as apresentações da Jornada de PCIs 2024, realizada em dezembro do ano passado.
Em 20 de fevereiro de 2025, convidamos os professores de Fatecs que apresentaram seus relatos de experiência na Jornada a produzirem textos para a seção “Artigo de Opinião”.
Como resultado, temos as contribuições do Núcleo de Estudos da Linguagem da Fatec (NELF) Bragança Paulista; da professora de inglês Talita Annunciato Rodrigues (Fatec Indaiatuba) e do professor de espanhol Odenildo França Almeida (Fatec Ipiranga).
O texto do NELF Bragança Paulista ressalta que o uso de diversos idiomas (português, inglês, espanhol) nos PCIs / Cesu contribui para a melhoria das competências linguísticas dos futuros tecnólogos, o que é imprescindível para a atuação profissional globalizada.
Talita Rodrigues descreve a experiência de aprendizado intercultural no PCI / Cesu na área de marketing internacional realizado entre a Fatec Indaiatuba e a University of Minnesota Crookston.
Odenildo França Almeida relata o projeto “Escrituras Narrativas”, desenvolvido entre Fatec Ipiranga e Uniminuto desde 2023.
Nas quatro edições, houve PCIs com cronogramas variando entre cinco semanas (“muito \textbf{corrido}”) e oito semanas (“duração ideal”, no entendimento do autor).
Odenildo destaca que a avaliação pós-projeto, feita pelos professores envolvidos, é “essencial para identificar pontos a serem ajustados e resgatar práticas bem-sucedidas”.
Também temos neste número, Selma de Fátima Grossi e Lígia Grandi que escrevem sobre a importância da sinergia entre professores de disciplinas profissionais (conteúdos técnicos) e de idiomas estrangeiros para o sucesso dos Projetos Colaborativos Internacionais (PCIs / Cesu).
Quer contribuir você também?
Acesse https://publicacoescesu.cps.sp.gov.br/VEm/about/submissions Dúvidas?
Escreva para cesu.pci@cps.sp.gov.br.
Boa leitura!

% matched lemmas: corrido
\end{textsample}
